 \documentclass[12pt,reqno]{amsart}
%%%%%%%%%%%%%%%%%%%%%%%%%%%%%%%%%%%%%%%%%%%%%%%%%%%
%								Packages									         %
%%%%%%%%%%%%%%%%%%%%%%%%%%%%%%%%%%%%%%%%%%%%%%%%%%%
\usepackage[T1]{fontenc}
\usepackage{amsmath, amssymb, amsthm, amscd, amsfonts}
\usepackage{mathtools}
\usepackage{enumerate,multicol}
\usepackage[shortlabels]{enumitem}

\usepackage[dvipsnames]{xcolor}
\definecolor{DartmouthGreen}{HTML}{00693E}
\definecolor{DartmouthDarkGreen}{HTML}{004B23}
\definecolor{DartmouthLightGreen}{HTML}{4BAF7A}

\usepackage[
  colorlinks=true,
  hyperindex,
  linkcolor=DartmouthDarkGreen,
  citecolor=DartmouthGreen,
  urlcolor=DartmouthGreen,
  pagebackref=true
]{hyperref}

\usepackage{parskip}
\usepackage[letterpaper, margin=.5in, top=1in, bottom=1.4in, headheight=42pt, headsep=14pt, footskip=50pt]{geometry}
\linespread{1.1}

\usepackage{algorithm,algpseudocode,verbatim}
\usepackage{float,caption,comment}

\usepackage{tikz,tikz-cd}
\usetikzlibrary{graphs, graphs.standard,positioning}



\usepackage{calrsfs}
\DeclareMathAlphabet{\pazocal}{OMS}{zplm}{m}{n}
\usepackage{newcent,fouriernc}

%%%%%%%%%%%%%%%%%%%%%%%%%%%%%%%%%%%%%%%%%%%%%%%%%%%
%								Copyright 								         %
%%%%%%%%%%%%%%%%%%%%%%%%%%%%%%%%%%%%%%%%%%%%%%%%%%%
\usepackage{ccicons}
\usepackage{fancyhdr}

\pagestyle{fancy}
\fancyhf{}
\fancyfoot[R]{\footnotesize \thepage}
\renewcommand{\headrulewidth}{0pt}
\renewcommand{\footrulewidth}{0pt}

\newcommand{\originalurl}{}
\newcommand{\originalurlI}{}

\fancypagestyle{firstpage}{%
  \fancyhf{}
\fancyhead[L]{\footnotesize Dartmouth College \\ Math 121: Commutative Algebra}
\fancyhead[R]{\footnotesize \href{https://www.juliettebruce.xyz/teaching/121S26}{Spring 2026}}
\fancyfoot[L]{\footnotesize \ccbyncsa\ \ \ccCopy \ Juliette Bruce 2026.\\
Licensed under \href{https://creativecommons.org/licenses/by-nc-sa/4.0/}{Creative Commons BY-NC-SA 4.0}.\\
Adapted from \ccCopy \href{https://sites.lsa.umich.edu/kesmith/}{Karen Smith} (2019); see \href{\originalurl}{original}.}
\fancyfoot[R]{\footnotesize \thepage}
\renewcommand{\headrulewidth}{0.4pt}
\renewcommand{\footrulewidth}{0.4pt}
}

\fancypagestyle{firstpage2}{%
  \fancyhf{}
\fancyhead[L]{\footnotesize Dartmouth College \\ Math 121: Commutative Algebra}
\fancyhead[R]{\footnotesize \href{https://www.juliettebruce.xyz/teaching/121S26}{Spring 2026}}
\fancyfoot[L]{\footnotesize \ccbyncsa\ \ \ccCopy \ Juliette Bruce 2026.\\
Licensed under \href{https://creativecommons.org/licenses/by-nc-sa/4.0/}{Creative Commons BY-NC-SA 4.0}.\\
Adapted from \ccCopy \href{https://sites.lsa.umich.edu/kesmith/}{Karen Smith} (2019); see \href{\originalurl}{original} and \href{\originalurlI}{original}.}
\fancyfoot[R]{\footnotesize \thepage}
\renewcommand{\headrulewidth}{0.4pt}
\renewcommand{\footrulewidth}{0.4pt}
}


\fancypagestyle{firstpageIP}{%
  \fancyhf{}
\fancyhead[L]{\footnotesize Dartmouth College \\ Math 121: Commutative Algebra}
\fancyhead[R]{\footnotesize \href{https://www.juliettebruce.xyz/teaching/121S26}{Spring 2026}}
\fancyfoot[L]{\footnotesize \ccbyncsa\ \ \ccCopy \ Juliette Bruce 2026.\\
Licensed under \href{https://creativecommons.org/licenses/by-nc-sa/4.0/}{Creative Commons BY-NC-SA 4.0}.\\
Inspired by \ccCopy\ \href{https://sites.lsa.umich.edu/kesmith/}{Karen Smith} (2019); \href{https://sites.lsa.umich.edu/kesmith/teaching/math-614-commutative-algebra/}{see}.}
\fancyfoot[R]{\footnotesize \thepage}
\renewcommand{\headrulewidth}{0.4pt}
\renewcommand{\footrulewidth}{0.4pt}
}
%%%%%%%%%%%%%%%%%%%%%%%%%%%%%%%%%%%%%%%%%%%%%%%%%%%
%								Theorems 								         %
%%%%%%%%%%%%%%%%%%%%%%%%%%%%%%%%%%%%%%%%%%%%%%%%%%%

\newtheorem{maintheorem}{Theorem}
\newtheorem{theorem}{Theorem}
\newtheorem{lemma}[theorem]{Lemma}

\newtheorem{corollary}[theorem]{Corollary}
\newtheorem{proposition}[theorem]{Proposition}
\newtheorem{conjecture}[theorem]{Conjecture}


\theoremstyle{definition}
\newtheorem{maindefinition}[maintheorem]{Definition}
\newtheorem{definition}[theorem]{Definition}
\newtheorem{setup}[theorem]{Set-Up}
\newtheorem{notation}[theorem]{Notation}
\newtheorem{convention}[theorem]{Convention}
\newtheorem{construction}[theorem]{Construction}





\setcounter{MaxMatrixCols}{18}

%%%%%%%%%%%%%%%%%%%%%%%%%%%%%%%%%%%%%%%%%%%%%%%%%%%
%								Letters									         %
%%%%%%%%%%%%%%%%%%%%%%%%%%%%%%%%%%%%%%%%%%%%%%%%%%%


\newcommand{\CC}{\mathbb{C}}
\newcommand{\FF}{\mathbb{F}}
\newcommand{\EE}{\mathbb{E}}

\newcommand{\GG}{\mathbb{G}}
\newcommand{\HH}{\mathbb{H}}
\newcommand{\II}{\mathbb{I}}
\newcommand{\KK}{\mathbb{K}}

\newcommand{\MM}{\mathbb{M}}
\newcommand{\NN}{\mathbb{N}}
\newcommand{\PP}{\mathbb{P}}
\newcommand{\QQ}{\mathbb{Q}}
\newcommand{\RR}{\mathbb{R}}
\renewcommand{\SS}{\mathbb{S}}
\newcommand{\TT}{\mathbb{T}}
\newcommand{\VV}{\mathbb{V}}
\newcommand{\WW}{\mathbb{W}}

\newcommand{\ZZ}{\mathbb{Z}}


\newcommand{\cA}{\mathcal{A}}
\newcommand{\cB}{\mathcal{B}}
\newcommand{\cC}{\mathcal{C}}
\newcommand{\cE}{\mathcal{E}}

\newcommand{\cF}{\mathcal{F}}
\newcommand{\cG}{\mathcal{G}}

\newcommand{\cH}{\mathcal{H}}
\newcommand{\cI}{\mathcal{I}}
\newcommand{\cL}{\mathcal{L}}
\newcommand{\cM}{\mathcal{M}}
\newcommand{\cO}{\mathcal{O}}
\newcommand{\cP}{\mathcal{P}}
\newcommand{\cQ}{\mathcal{Q}}
\newcommand{\cS}{\mathcal{S}}
\newcommand{\cT}{\mathcal{T}}
\newcommand{\cR}{\mathcal{R}}
\newcommand{\cU}{\mathcal{U}}
\newcommand{\cV}{\mathcal{V}}

\newcommand{\pB}{\pazocal{B}}
\newcommand{\pC}{\pazocal{C}}
\newcommand{\pO}{\pazocal{O}}
\newcommand{\pG}{\pazocal{G}}

\newcommand{\sM}{\mathsf{M}}
\newcommand{\sN}{\mathsf{N}}

\newcommand{\sQ}{\mathsf{Q}}
\newcommand{\sP}{\mathsf{P}}

\newcommand{\fm}{\mathfrak{m}}
\newcommand{\fp}{\mathfrak{p}}
\newcommand{\fq}{\mathfrak{q}}

%%%%%%%%%%%%%%%%%%%%%%%%%%%%%%%%%%%%%%%%%%%%%%%%%%%
%								Operators									         %
%%%%%%%%%%%%%%%%%%%%%%%%%%%%%%%%%%%%%%%%%%%%%%%%%%%
\DeclareMathOperator{\Id}{Id}
\DeclareMathOperator{\ev}{ev}
\DeclareMathOperator{\ord}{ord}
%
\DeclareMathOperator{\Hom}{Hom}
\DeclareMathOperator{\End}{End}
\DeclareMathOperator{\coker}{coker}
\DeclareMathOperator{\Ann}{Ann}
\DeclareMathOperator{\img}{img}
%
\DeclareMathOperator{\Frac}{Frac}
\DeclareMathOperator{\Nil}{Nil}
\DeclareMathOperator{\red}{red}
%
\DeclareMathOperator{\Spec}{Spec}
\DeclareMathOperator{\mSpec}{mSpec}
%
\DeclareMathOperator{\SL}{SL}
%
\DeclareMathOperator{\init}{in}
\DeclareMathOperator{\lex}{lex}
\DeclareMathOperator{\grlex}{grlex}
\DeclareMathOperator{\grevlex}{grevlex}
\DeclareMathOperator{\revlex}{revlex}
\DeclareMathOperator{\LT}{LT}
\DeclareMathOperator{\LM}{LM}
\DeclareMathOperator{\LC}{LC}


%%%%%%%%%%%%%%%%%%%%%%%%%%%%%%%%%%%%%%%%%%%%%%%%%%%
%								Categories									         %
%%%%%%%%%%%%%%%%%%%%%%%%%%%%%%%%%%%%%%%%%%%%%%%%%%%

\DeclareMathOperator{\comRing}{\textbf{comRing}}
\DeclareMathOperator{\Top}{\textbf{Top}}
\DeclareMathOperator{\Set}{\textbf{Set}}
\DeclareMathOperator{\alg}{\textbf{alg}}
\DeclareMathOperator{\Mod}{\textbf{Mod}}


%%%%%%%%%%%%%%%%%%%%%%%%%%%%%%%%%%%%%%%%%%%%%%%%%%%
%								MISC									         %
%%%%%%%%%%%%%%%%%%%%%%%%%%%%%%%%%%%%%%%%%%%%%%%%%%%
\makeatletter
\newcommand*{\tp}{%
  {\mathpalette\@transpose{}}%
}
\newcommand*{\@transpose}[2]{%
  % #1: math style
  % #2: unused
  \raisebox{\depth}{$\m@th#1\intercal$}%
}
\makeatother

\newcommand{\juliette}[1]{{\color{purple} \sf  Juliette: [#1]}}
%%%%%%%%%%%%%%%%%%%%%%%%%%%%%%%%%%%%%%%%%%%%%%%%%%%
%%%%%%%%%%%%%%%%%%%%%%%%%%%%%%%%%%%%%%%%%%%%%%%%%%%
%%%%%%%%%%%%%%%%%%%%%%%%%%%%%%%%%%%%%%%%%%%%%%%%%%%
%%%%%%%%%%%%%%%%%%%%%%%%%%%%%%%%%%%%%%%%%%%%%%%%%%%
\renewcommand{\originalurl}{https://sites.lsa.umich.edu/kesmith/wp-content/uploads/sites/1309/2024/07/ValuationRings.pdf}
\title{Worksheet 2.3: Valuation Rings}

%%%%%%%%%%%%%%%%%%%%%%%%%%%%%%%%%%%%%%%%%%%%%%%%%%%
%%%%%%%%%%%%%%%%%%%%%%%%%%%%%%%%%%%%%%%%%%%%%%%%%%%

%%%%%%%%%%%%%%%%%%%%%%%%%%%%%%%%%%%%%%%%%%%%%%%%%%%
%%%%%%%%%%%%%%%%%%%%%%%%%%%%%%%%%%%%%%%%%%%%%%%%%%%
\begin{document}

%%%%%%%%%%%%%%%%%%%%%%%%%%%%%%%%%%%%%%%%%%%%%%%%%%%
%%%%%%%%%%%%%%%%%%%%%%%%%%%%%%%%%%%%%%%%%%%%%%%%%%%

\maketitle
\thispagestyle{firstpage}

Fix a field $\KK$. In this worksheet we will explore a special, but important type of rings called valuation rings. 
\begin{definition}
	A \emph{valuation ring} of a field $\KK$ is a subring $R \subset \KK$ with the property that for all $x \in \KK^{\times}$ either $x \in R$ or $x^{-1} \in R$. 
\end{definition}

Note that, since a valuation ring $R$ of a field $\KK$ is a subring of $\KK$ it is necessarily a domain. We first study valuation rings directly, then show how certain functions that broadly measure some notion of divisibility, called valuations, produce them. Lastly, we compare discrete and non-discrete examples.


 \hrulefill 

\begin{enumerate}[leftmargin=*]
\item \textbf{Examples of Valuation Rings.} Recall $\CC(x,y)$ is the field of formal rational functions in $x$ and $y$, which concretely is the set of elements $f/g$ where $f,g\in \CC[x,y]$ with $g\neq 0$.  For each of the following subsets $\cS \subset \CC(x,y)$ determine whether or not they are valuation rings by showing whether for every $h \in \CC(x,y)^{\times}$ either $h \in \cS$ or $h^{-1} \in \cS$.
\begin{enumerate}[label=(\alph*),ref=\theenumi\alph*]
    \item Hint: Use that $\CC[x,y]$ is a UFD to understand when $g$ is divisible by $y$.
    \[
    \cS_{1} \coloneqq \left\{ \frac{f}{g} \;\; \big| \;\; f,g \in \CC[x,y] \text{ and } g\not \in \langle y \rangle \right\}
    \]


    \item Hint: Consider the element $\frac{x}{y}$. 
    \[
    \cS_{2} \coloneqq \left\{ \frac{f}{g} \;\; \big| \;\; f,g \in \CC[x,y] \text{ and } g(0,0) \neq 0 \right\}
    \]

    \item Define the \emph{order} of a non-zero polynomial $f \in \mathbb{C}[x,y]$ to be $\ord(f)$, the lowest degree of a monomial in $f$ and $\ord(0)=\infty$. 
    \[
    \cS_{3} \coloneqq \left\{ \frac{f}{g} \;\; \big| \;\; f,g \in \CC[x,y], \; g\neq0, \text{ and } \ord(f) \geq \ord(g) \right\}
    \]
\end{enumerate}

 \item \textbf{Ideals in Valuation Rings.} Let $R$ be a valuation ring of $\KK$. 
      \begin{enumerate}[label=(\alph*),ref=\theenumi\alph*]
      \item Show that for any two non-zero elements $x, y \in R$, either $x|y$ or $y|x$.
      \item Show that every finitely generated ideal of  $R$  is principal.
      \item Prove that the ideals of $R$ form a totally ordered set. That is, for any two ideals $I, J$ in $R$, either $I \subseteq J$ or $J\subseteq I$. (Hint: Given $x\in I$ with $x\not\in J$ consider whether an element $y\in J$ can divide $x$.)
      \item What does the previous parts say about the poset structure on $\Spec(R)$?
\end{enumerate}

\item \textbf{Local Rings \& Valuation Rings.} An arbitrary ring $R$ is said to be a \emph{local ring} if it has exactly one maximal ideal. We often write a local ring $R$ as a pair $(R,\fm)$ where $\fm \subset R$ is the unique maximal ideal. 
\begin{enumerate}[label=(\alph*),ref=\theenumi\alph*]
	\item Describe $\mSpec(R)$ if $R$ is a local ring.
	\item Describe the closed points of $\Spec(R)$ if $R$ is a local ring. 
	\item Prove that a valuation ring is a local ring. 
\end{enumerate}

\item \textbf{$p$-Adic Integers.} Fix a prime $p$. The \emph{$p$-adic numbers} $\mathbb{Q}_p$ are the completion of $\mathbb{Q}$ with respect to the $p$-adic absolute value $|\cdot|_p$, defined as follows: for a nonzero rational number $x$, write $x = p^k \frac{a}{b}$ where $a,b \in \mathbb{Z}$ are not divisible by $p$. Then
\[
|x|_p = p^{-k}, \qquad |0|_p = 0.
\]
Note $|\cdot|_{p}$ extends uniquely to $\QQ_{p}$. The \emph{$p$-adic integers} are defined by
\[
\ZZ_p = \left\{x \in \QQ_p \;\; \big| \;\;  |x|_p \leq 1\right\}.
\]
Similar to decimal expansion in $\RR$, it is a fact that every element $x \in \QQ_p$ can be written uniquely as a formal series
\[
x = \sum_{i=k}^{\infty} a_i p^i  \quad \quad \text{where} \quad \quad a_i \in \left\{0,1,\dots,p-1\right\}, \quad a_k \neq 0,
\]
and that
\[
\ZZ_p = \left\{ \sum_{i=0}^{\infty} a_i p^i \;\; \big| \;\; a_i \in \{0,1,\dots,p-1\} \right\}.
\]
\begin{enumerate}[label=(\alph*),ref=\theenumi\alph*]
   \item Prove that the $p$-adic absolute value is well-defined. (Hint: Use the fact that $\ZZ$ is a UFD to show that every element $x\in \QQ$ can be written uniquely (up to sign and reordering) as  $p_1^{a_1}\dots p_t^{a_t}$ where the $p_i$ are distinct primes and $a_i \in \mathbb Z$.)
    \item\label{ex-pt:padic-add} Prove the $p$-adic absolute value satisfies the \emph{non-Archimedean (ultrametric) inequality}:
    \[
    |x+y|_p \leq \max\left\{|x|_p, |y|_p\right \}.
    \]
    
    \item\label{ex-pt:padic-mlt}  Prove the $p$-adic absolute value is multiplicative: $|xy|_{p} = |x|_{p}|y|_{p}$.
    
    \item Use parts \ref{ex-pt:padic-add} and \ref{ex-pt:padic-mlt} to show that $\ZZ_p$ is closed under addition and multiplication.
    
    \item Conclude that $\ZZ_p$ is a ring (called the ring of $p$-adic integers).
    
    \item Prove that $\ZZ_{p}$ is a valuation ring. 
\end{enumerate}
\end{enumerate}

\hrulefill

The previous exercises suggest why valuation rings are called \emph{valuation} rings: they arise from functions that measure divisibility, generalizing the $p$-adic absolute value. Such functions are called valuations. An \emph{ordered abelian group} is an abelian group $(\Gamma,+)$ together with a total order $<$ on $\Gamma$ that is compatible with addition, in the sense that if $\alpha<\beta$ then  $\alpha+\gamma<\beta+\gamma$ for all $\alpha,\beta,\gamma \in \Gamma$.   We will frequently extend the total ordering on an ordered abelian group $\Gamma$ to $\Gamma \cup \{\infty\}$ by making the convention that $\gamma < \infty$ for all $\gamma \in \Gamma$. Similarly, we will extend addition to $\Gamma\cup\{\infty\}$ by declaring $\infty + \gamma = \gamma + \infty = \infty$ for all $\gamma \in \Gamma$.

\begin{definition}
Let $(\Gamma,+,<)$ be an ordered abelian group. A  \emph{$\Gamma$-valued valuation}, or just a \emph{valuation}, on $\KK$ is a function: $\nu:\KK \to \Gamma\cup\{\infty\}$, which satisfies the following axioms for all $x,y\in \KK$:
\begin{enumerate}[(i)]
\item $\nu(xy) = \nu(x) + \nu(y)$
\item $\nu(x + y) \geq \min\left\{ \nu(x), \nu(y)\right\}$, and
\item $\nu(x) = \infty$ if and only if $x=0$.
\end{enumerate}
\end{definition}

Note that this means $\nu$ defines a group homomorphism $\nu:\KK^{\times} \to \Gamma$. Indeed, it is standard to define a valuation as a group homomorphism $\nu: \KK^{\times} \to \Gamma$ satisfying axiom (ii), and then extend it to all of $\KK$ by setting $\nu(0)=\infty$. This convention is implicit throughout the remainder of the worksheet. 

The \emph{value group} of a valuation $\nu:\KK^{\times} \to \Gamma$ is the image of $\nu$. A valuation is a \emph{discrete valuation} if its value group is isomorphic to $\ZZ$. If $\FF \subset \KK$ is a subfield, we say that $\nu$ is a $\FF$-valuation if $\nu(\lambda) = 0 $ for all $\lambda \in \FF^{\times}$.

\hrulefill

\begin{enumerate}[leftmargin=*]
  \setcounter{enumi}{4}

\item \textbf{Examples of Valuations.} Prove that the following are valuations.
	\begin{enumerate}[label=(\alph*),ref=\theenumi\alph*]
	\item The $p$-adic valuation:
	\[
	\nu_p:\QQ^\times \to \ZZ,
	\qquad
	\nu_p\!\left(p^k\frac{a}{b}\right)=k
	\]
	where $a,b\in \ZZ$ are not divisible by $p$.
	
	\item The order valuation on $\CC(x)$:
	\[
	\nu_0:\CC(x)^\times \to \ZZ,
	\qquad
	\nu_0(f/g)=\ord_{x=0}(f)-\ord_{x=0}(g).
	\]
	
	\item More generally, if $a\in \CC$, define
	\[
	\nu_a:\CC(x)^\times \to \ZZ
	\]
	by order of vanishing at $x=a$. Show this is again a valuation.
	
	\item On $\CC(x,y)$, define
	\[
	\nu(f/g)=\operatorname{ord}(f)-\operatorname{ord}(g),
	\]
	where $\operatorname{ord}(h)$ is the degree of the lowest-degree nonzero term of $h\in \CC[x,y]$. Show that this is a valuation with value group $\ZZ$.
	\end{enumerate}
	
\item Prove that every ordered abelian group is torsion-free. (Hint: Reduce to the case when $x>0$.)

\item \textbf{Basic Properties of Valuations.} Let $\nu:\KK^{\times} \to \Gamma$ be an arbitrary valuation. Fix $f,g \in \KK^{\times}$.
\begin{enumerate}[label=(\alph*),ref=\theenumi\alph*]
	\item Show that $\nu(f^{-1}) = -\nu(f)$.
	\item\label{ex-pt:val-sum} Show that if $\nu(g) < \nu(f)$, then $\nu(f+g) = \nu(g).$ (Hint: Use that $(f+g) + (-f) = g$.)
	\item Show that  $\nu(f) = \nu (-f)$. (Hint: Use part \ref{ex-pt:val-sum}.)
\end{enumerate}


 \item Consider the field $ \CC(x, y)$. Throughout this question $\pi$ refers to the number $\pi=3.1415926...$.
     \begin{enumerate}[label=(\alph*),ref=\theenumi\alph*]
      \item\label{ex-pt:val-exm1} Show that there is a unique $\CC$-valuation $\nu_{\pi}:\CC(x,y)^{\times} \to (\RR,+)$  such that 
      \[
      \nu_{\pi}\left(x^ay^b\right) = a + \pi b.
      \]
      What is the value group of $\nu_{\pi}$? (Hint: Use Question 3.) 
      \item\label{ex-pt:val-exm2} Consider the group $\Gamma = \ZZ\oplus\ZZ$ ordered by the lexicographical ordering. Show that there exists a unique $\CC$-valuation $\nu:\CC(x,y)^{\times} \to \Gamma$ such that $\nu(x^ay^b)=(a,b)$. 
      \item Are either of these valuations constructed in parts \ref{ex-pt:val-exm1} or \ref{ex-pt:val-exm2} discrete?
      \end{enumerate}
        
 
 \item \textbf{Valuations Determine Valuation Rings.}  Let $\nu$ be an arbitrary valuation on $\KK$. 
	\begin{enumerate}[label=(\alph*),ref=\theenumi\alph*]
		\item Prove that the set $R \subset \KK$ of elements $f\in \KK$ with $\nu(f) \geq 0$ is a valuation ring of $\KK$. We call this the \emph{valuation ring} of $\nu$. (Hint: Remember our convention about $\nu(0)$.)
		\item Show that the following is an ideal of $R$:
		\[
		\fm \coloneqq \left\{ f \in \KK \;\; \big| \;\; 0 < \nu(f) \right\}.
		\]
		\item Prove that if $f \in R$ and $f \not \in \fm$ then $f$ is a unit.
		\item Combine the previous parts to deduce $\fm$ is the unique maximal ideal of $R$. 
		\item Fix an element $\gamma \in \Gamma$. Prove the sets below are ideals of $R$:
		\[
		\left\{ f \in R \;\; \big|\;\; \nu(f) > \gamma \right\} \quad \text{and} \quad \left\{ f \in R \;\; \big|\;\; \nu(f) \geq \gamma \right\}.
		\]
	\end{enumerate}
 
\item \textbf{Discrete Valuation Rings.} Let $\nu:\KK^{\times} \to \ZZ$ be a discrete valuation on $\KK$. Let $(R,\fm)$ be the corresponding \emph{discrete valuation ring}. 
 \begin{enumerate}[label=(\alph*),ref=\theenumi\alph*]
	\item Let $I \subset R$ be any non-zero ideal, explain why we can choose an element $f \in I $ such that $\nu(f)$ is minimal among elements in $I$.
	\item\label{ex-pt:dvr-pid}  Prove that every ideal in $R$ is principal. In particular, a discrete valuation ring is a PID. 
	\item By part \ref{ex-pt:dvr-pid} there exists an element $t\in R$ such that $\langle t \rangle = \fm$. Show that every ideal of $R$ is either $0$ or of the form $\langle t^{n} \rangle$ for some $n\geq0$. 
\end{enumerate}

\item Consider the valuation $\nu_{\pi}$ on $\CC(x,y)$ constructed in Question~\ref{ex-pt:val-exm1}. Let $(R,\fm)$ be the corresponding valuation ring for $\nu_{\pi}$.
 \begin{enumerate}[label=(\alph*),ref=\theenumi\alph*]
 	\item For a real number $x \in \RR_{\geq 0}$ let $\{x\}\coloneqq x - \lfloor x \rfloor$ denote its fractional part. Fix a positive integer $N$ and consider the sequence consisting of the fractional parts of the first $N+1$ multiples of $\pi$: $\{0\}, \{\pi\}, \{2\pi\},\ldots,\{N\pi\} \in [0,1)$. Divide the interval $[0,1)$ into $N$ evenly spaced subintervals: 
	\[
	\left[0, 1\right) = \left[0,\frac{1}{N}\right) \cup \left[\frac{1}{N},\frac{2}{N}\right)\cup \cdots \cup \left[\frac{N-1}{N},1\right).
	\] 
	Prove there exist integers $0\leq k,j\leq N$ such that $\{k\pi\}$ and $\{j\pi\}$ lie in the same subinterval. Deduce that for any positive integer $N$ there exist integers $0\leq k,j\leq N$ such that $|\{k\pi\} - \{j\pi\}| \leq \frac{1}{N}$. (Hint: Think about pigeons.)
	\item Using that $\pi$ is irrational show that $\{j\pi\}\neq \{k\pi\}$ whenever $j\neq k$. 
 	\item\label{ex-pt:ZplusPiZ-dense} Prove for every real number $\epsilon>0$, there exist integers $a,b\in \ZZ$ such that $0 < a+b\pi < \epsilon$. 
	\item Prove that $\fm$ is not finitely generated. (Hint: Use part \ref{ex-pt:ZplusPiZ-dense}.)
	\item Prove that $\fm$ is the only non-zero radical ideal of $R$. (Hint: If $I \subset R$ is a non-zero radical ideal, let $f \in I$ be non-zero. For an arbitrary element $h \in \fm$ show that $h^{n} \in \langle f \rangle$ for $n \gg0$.)
	\item Describe $\Spec(V)$. 
\end{enumerate}
 
 \item Consider the valuation $\nu$ on $\CC(x,y)$ constructed in Question~\ref{ex-pt:val-exm2}. Let $(R,\fm)$ be the corresponding valuation ring for $\nu$.
  \begin{enumerate}[label=(\alph*),ref=\theenumi\alph*]
 	\item Show that the maximal ideal of $R$ is
	\[
	\fm = \left\{f \in \CC(x,y) \;\; \big| \;\; \nu(f) \geq  (0, 1)\right\} = \langle y \rangle.
	\]
 	\item Show that the ideal  $\langle x\rangle \subset R$ is not prime. (Hint: $x = xy^{-1}\cdot y$.)
	\item Show that the set below is a prime, non-maximal ideal of $R$:
	\[
	P \coloneqq \left \{ f \in R \;\; \big | \;\; \nu(f) = (a,b) \text{ where } a \geq 1 \right\}\cup\{0\}.
	\] 
 	\item Use the valuation to construct an infinite strictly ascending chain of ideals in $R$:
	\[
	\left\langle xy^{-1} \right\rangle \subset \left\langle xy^{-2} \right\rangle \subset \left\langle xy^{-3}  \right\rangle \subset \cdots.
	\]
	Explain why all inclusions are strict.
  \end{enumerate}

     
\item \textbf{Valuation Rings Determine Valuations.}
Let $R$ be a valuation ring of $\KK$ with $R \neq \KK$. This problem shows how to recover a valuation from the divisibility structure of $R$.

\begin{enumerate}[label=(\alph*),ref=\theenumi\alph*]
	\item Let $\Gamma'$ be the set of nonzero principal ideals of $R$:
	\[
	\Gamma' = \left\{\langle a\rangle \subset R \;\; \big|\;\; a \in R,\ a \neq 0\right\}.
	\]
	Define a multiplication on $\Gamma'$ by ideal multiplication: $\langle a\rangle \cdot \langle b\rangle = \langle ab\rangle$.
	\begin{enumerate}
		\item Show that this operation is well-defined and makes $\Gamma'$ into an abelian monoid (associative, commutative, binary operation with an identity element).
		\item Show that $\Gamma'$ is totally ordered by inclusion: for distinct elements $\langle a\rangle,\langle b\rangle \in \Gamma'$, either $\langle a\rangle \subseteq \langle b\rangle$ or $\langle b\rangle \subseteq \langle a\rangle$.
		\end{enumerate}

	\item Define an equivalence relation $\sim$ on $\KK^\times$ by $x \sim y$ if and only if $x=uy$ for some unit $u \in R^{\times}$. Let $\Gamma = \KK^\times / \sim$ denote the set of equivalence classes. Given $x\in \KK^{\times}$ write $[x]$ its equivalence class in $\Gamma$.
	\begin{enumerate}
		\item Show that $\sim$ is a well-defined equivalence relation.
		\item Show that multiplication in $\KK^\times$ induces a well-defined operation on $\Gamma$, making it into an abelian group.
		\item Define $[x] \leq [y]$ if and only if $yx^{-1} \in R$. Show this is a well-defined total order on $\Gamma$.
		\item Deduce that $\Gamma$ is an ordered abelian group.
		\item Show that the map $\Gamma' \to \Gamma$ given by $\langle a \rangle \mapsto [a]$ is an order-preserving embedding whose image consists exactly of the elements $\ge [1]$.
	\end{enumerate}
	
	\item Define $\nu:\KK^{\times} \to \Gamma$ by $\nu(x) = [x]$. Prove that $\nu$ is a valuation on $\KK$ with value group $\Gamma$ and 
	\[
	R = \left\{ x \in \KK \;\; \big| \;\; \nu(x) \geq [1] \right\}.
	\]
	Moreover, $\nu(x) \in \Gamma'$ if and only if $x \in R$. 
	\end{enumerate}
  
\end{enumerate}

\end{document}





